\section{从缉私到条约口岸：行政选择如何失去余地}

鸦片问题并不是一个突然落到北京的道德难题。白银流动、海防责任、行商债务和沿海走私已让地方行政承压；1830年代朝廷内部仍有弛禁、征税、严缉等相互竞争的方案。1838年林则徐赴广东，次年虎门销烟，把原来在商人、海关和巡防之间反复推诿的责任集中为国家行动。这个选择有其财政和主权逻辑，也使清廷进入与英国武力和外交资源不对称的冲突。

1840年战争爆发，1842年《南京条约》签订，补充条约和五口通商随后才逐步进入地方行政。它们改变了关税、城市管理、领事和贸易的制度条件，却不能把签字那一刻写成全国秩序的瞬时转换。广州的入城争议、传教活动引发的地方案件以及长江沿岸的新商路，都显示条约权利必须在城墙、衙门、行会、教堂与社区的具体冲突中才有实际形状。

对普通家庭而言，危机并不只发生在外交桌前。河工失修、漕粮转运、旗人债务、灾害赈济和矿区迁徙继续决定谁能取得粮食、信用和工作；口岸的变化又把这些问题同银价、航运和新税源接在一起。清廷并非完全没有政策选择，而是每一种选择都要穿过机构分工、地方利益与资源不足，因而难以同时解决海防、财政和社会秩序。

\subsection{材料与争议}

《筹办夷务始末》、实录、照会和条约文本各自记录谈判与国家主张。中外文本必须区分签署、批准、换约和执行；行商账簿、地方档案、报刊与外国商务记录只能检验某些口岸。走私规模、银价效应和居民态度不能由任何一套档案单独裁定。

\subsection{交叉阅读窗口：魏源把“夷”变成一门必须学习的技术}

魏源在1843年初刻、1852年增订的《海国图志》中提出“师夷长技以制夷”。这是一部在战败后汇集地理、军政与海外知识的经世著作，不是1839年决策现场的实录；“师”与“制”的对置却把承认差距转为自强的动作，因而有很强的论辩与安定人心的力量。它应与《道光实录》、林则徐奏折和《筹办夷务始末》并读，以分清禁烟、战事和议和在清廷文书链中的先后；再把《南京条约》的中英文文本、广州行商账簿、地方档案、英国外交部档案与口岸英文报刊并置，才可追问条款如何在具体港口被翻译、争执和执行。

钱穆会把问题放在旧有财政、海防与外交机构能否吸收新知识；吕思勉则会问白银、商路、行商债务和沿海生计怎样构成这场危机的社会条件。魏源的文辞可以证明一位士人如何构想出路，不能证明朝廷采纳、技术输入成功或民众已经同意。条约能证明签署与承诺，外国公文和报刊则分别服务谈判、商业与侨民读者；它们都不能单独替广州、上海或内地居民发言。
